\itemhead{p.47 练习 · 第2题（教材p.47）}{XJ-B1-C02-PR-P047-Y1307}
\label{XJ-B1-C02-PR-P047-Y1307}
{\footnotesize 湘教版教材p.47 / PDF第54页 / p.47 练习 / 第2题}
\field{原题}{已知函数 $y=-x^2+2x-m+1$。
\subq{1}若 $x=0$ 为函数的一个零点，求 $m$ 的值；
\subq{2}当 $m$ 为何值时，函数有两个零点、一个零点、无零点？}
\field{参考答案}{（1）$m=1$；（2）由 $\Delta=8-4m$，$m<2$ 时两个零点，$m=2$ 时一个零点，$m>2$ 时无零点。}
{\small \field{关联知识}{Q01 二次函数、零点、判别式与根系关系（教材帮p.65 / PDF74）}
\field{方法与易错点}{按已知零点、顶点或系数选择交点式、顶点式或一般式；根平方和用根系关系消元。 易错：遗漏 $a
e0$；把二重根算成两个不同零点；有两正根只查积正。}
\field{答案依据}{答案册相关题组：PDF第13页（册内页12）；按当前题库原题逐项重算，不按旧版题号机械对套。}}
